\section{Intro}

With the background material covered, we can start getting to the good stuff.
%Overall, this chapter is a mixture of stuff that one can find already in the literature (but perhaps not all in one place), and new results that are essentially just me digging a little deeper into the basic material behind dynamic stabilizer codes.
%We begin with a thorough introduction to said codes.
In this chapter, we give a thorough introduction to dynamic stabilizer codes.
As mentioned in \Cref{ch:intro}, these first appeared in Hastings' and Haah's paper \Ccite{Hastings_2021} in 2021.

We begin by diving deeper into the stabilizer formalism, since we'll need a much better understanding of it in order to introduce dynamic stabilizer codes.
Much of the content in \Cref{sec:more-stabilizer-formalism} is standard, if underrepresented in the literature -- in particular, the main source for this section is Gottesman's lecture notes \Ccite{gottesman2024surviving}.
However, I'm always a little unsatisfied by any existing presentation I can find of this content.
Primarily, my gripe is that in the few sources that discuss how $\NNN(\SSS)/\SSS$ changes under the certain operations (Cliffords and Pauli measurements), this is phrased in terms of simple update rules -- i.e.\ how to present the new group $\NNN(\SSS')/\SSS'$ given a presentation for the old one $\NNN(\SSS)/\SSS$.
But this misses the most interesting and important part -- namely which logical operators in $\NNN(\SSS)/\SSS$ before the operation correspond to which logical operators in $\NNN(\SSS')/\SSS'$ after the operation!

Indeed, this perspective is key to understanding what I consider to be the most interesting feature of dynamic stabilizer codes -- that just by dint of the sequence of Cliffords and Pauli measurements that defines them (their \textit{operation schedule}), they can `naturally' perform logical Clifford operations.
For example, take the honeycomb code defined by Hastings and Haah in the paper that introduced dynamic stabilizer codes~\cite{Hastings_2021}.
When implemented on a torus, this encodes two qubits, and every three timesteps performs the logical gate $\pm (H \otimes H) \SWAP$.
Or better still, consider Davydova et al.'s scheme for implementing the entire multi-qubit logical Clifford group just by changing the operation schedule of a certain dynamic stabilizer code (which one can consider to be a \textit{Floquet colour code} -- more on this later)~\cite{Davydova_2024}.
We therefore take the time to spell this relationship out; for any Clifford unitary or Pauli measurement, we define an \textit{induced homomorphism} $\phi: \NNN(\SSS')/\SSS' \to \NNN(\SSS)/\SSS$ that tracks how the logical operators before and after the operation are related to one another.

We then give a definition of a general dynamic stabilizer code.
As a community, it seems we still haven't settled on any one definition for a dynamic stabilizer code.
The story from my perspective is something like this -- initially, Hastings and Haah's original paper just gave a couple of examples of dynamic stabilizer codes, and no general definition.
Christophe Vuillot's paper \Ccite{vuillot2021planarfloquetcodes} then introduced the name \textit{Floquet code} for a specific type of dynamic stabilizer code derivable from the colour code, of which Hasting's and Haah's code is one example.
This family was given other names in other places -- dynamic tree codes~\cite{Davydova_2023}, dynamically condensed colour codes~\cite{Kesselring_2024}, etc.
But then, over the next few months (now years), `Floquet code' seemed to come to mean any code defined by a schedule of Pauli measurements and Clifford operations~\cite{bauer2025xyfloquetcodesimple}.
While suitably general, this had the drawback that `Floquet' in physics apparently usually means `periodic', so it was perhaps a little confusing to view codes with infinite schedules~\cite{gidney2023baconthreshold} as Floquet codes.

Therefore, to be as general as possible while remaining within the stabilizer formalism framework, we sidestep the word `Floquet' and define a \textit{dynamic stabilizer code} to be any sequence (perhaps infinite) of Pauli measurements and Clifford unitaries.
In other words, any circuit consisting of Cliffords and Pauli measurements is a QEC code!
This is the mantra we mentioned in \Cref{ch:intro}, and has been one of the major themes in QEC over the last few years; circuits have been upgraded from objects that \textit{implement} codes to objects that \textit{are} codes.

Having introduced dynamic codes algebraically, we then take a look at them graphically via \textit{Pauli webs}.
Pauli webs first appeared independently in two papers, one from PsiQuantum~\cite{Bombin_2024} and one from Google~\cite{McEwen_2023}.
As mentioned in the introduction to this thesis, one can think of Pauli webs as graphical stabilizer formalism, which explains why they come up when discussing dynamic stabilizer codes.
For me personally, they turn working with dynamic stabilizer codes from something fiddly and offputting to something more tractable and rather fun.
In \Cref{sec:pauli-webs} we therefore give a thorough first principles introduction to them, then connect them to the stabilizer formalism and dynamic stabilizer codes.
Throughout this section, we often reference Rodatz et al.'s excellent papers~\cite{rodatz2024floquetifying,rodatz2025faulttoleranceconstruction}.
